% GDZ PPN236018647,native550upper;HKUST551.
\clearpage\SourceLeaf{550}\ArithmeticHead{521}{R}{1797 IUL. 23—1797 AUG. 1.}
\NumberedItem{71}
\indent Radicum aeq[uationis] \(x^n=1\) periodi plures eandem summam habere non possunt demonstratur.
\par\noindent\hfill[1797] Iul. 27.\quad Gott[ingae]

\medskip\begingroup\small
\indent Vergl. die Notiz Nr. 73 vom 1. August 1797.
\par\endgroup

\NumberedItem{72}
\indent Plani possibilitatem demonstravi.
\par\noindent\hfill[1797] Iul. 28.\quad Gotting[ae]

\medskip\begingroup\small
\indent Am 6. März 1832 schreibt \textsc{Gauss} an \textsc{Wolfgang Bolyai} (Werke VIII, S. 224): »Um die Geometrie vom Anfange an ordentlich zu behandeln, ist unerläßlich, die Möglichkeit eines Planums zu beweisen; die gewöhnliche Definition enthält zu viel, und implicirt eigentlich subreptive schon ein Theorem. Man muß sich wundern, daß alle Schriftsteller von \textsc{Euklid} bis auf die neuesten Zeiten so nachlässig dabei zu Werk gegangen sind«. Wohl veranlaßt durch den Brief an \textsc{Bolyai} hat \textsc{Gauss} im März 1832 seine \textit{Begründung des Planum} im Handbuch 19, Be, S. 153 aufgezeichnet, sie ist abgedruckt Werke VIII, S. 194; an derselben Stelle findet man auch unter [2.] eine kurze, aus der Zeit zwischen 1820 und 1830 stammende Notiz über die Notwendigkeit eines Beweises für die Möglichkeit der Ebene. Man vergleiche ferner den Brief an \textsc{Bessel} vom 27. Jan. 1829, Werke VIII, S. 200.
\par\noindent\hfill\textsc{Stäckel.}
\par\endgroup
